\section{Thermal transfer construction and stable realization}
\label{sec:thermal-realization}

\subsection{Source and observation operators}

Let the electromagnetic loss model provide $R_q$ source profiles
$\Psi_\ell(\mathbf r)$ with amplitudes $p_\ell(t)$. Define
\[
 q(\mathbf r,t)
 =
 \sum_{\ell=1}^{R_q}\Psi_\ell(\mathbf r)p_\ell(t).
\]
Let $\mathcal O_T$ be a set of thermal observations: component-average
temperature, point probes, surface temperatures, or coefficients of a
continuous decoder basis.

After Laplace transformation, the thermal teacher defines
\[
 y_T(s)=H_T(s;\mathcal S)p(s)+h_{\rm env}(s),
\]
where
\[
 H_T(s)
 =
 C_{\rm obs}
 \mathcal L_T(s)^{-1}
 B_{\rm src}.
\]
$\mathcal L_T(s)$ is the continuous conduction/interface operator.

\subsection{Sampling in the Laplace domain}

The transfer is sampled at shifts
\[
 s_k=\sigma_k+j\omega_k,\qquad \sigma_k>0
\]
covering the time scales of interest. Sampling points are selected
logarithmically and adaptively according to transfer interpolation error. Very
long-time behavior is constrained by a separate near-zero/steady solve rather
than extrapolated from only short-time samples.

\subsection{Structure-preserving reduction}

A linear passive thermal transfer is a positive-real/Stieltjes-type diffusive
system under appropriate source/observation pairing. The reduced realization
is sought directly in the form
\[
 C_T\dot z=-G_Tz+B_Tp,\qquad
 y=C_T^{\rm obs}z+D_Tp,
\]
with
\[
 C_T\SPD,\qquad G_T\PSD.
\]
Projection bases obtained from rational Krylov or moment matching may be used,
but the final realization is symmetrized/parameterized so the energy
\[
 \mathcal E_T(z)=\frac12 z^\T C_Tz
\]
satisfies
\[
 \frac{\dd\mathcal E_T}{\dd t}
 =
 -z^\T G_Tz+z^\T B_Tp
\]
up to environmental exchange terms.

\subsection{Geometry-parametric thermal model}

The thermal state dimension is fixed only within one artifact schema, not by a
world mesh. Geometry dependence enters through structured factors:
\[
 C_T(\mathcal S)=L_C(\mathcal S)L_C(\mathcal S)^\T+\epsilon I,
\]
\[
 G_T(\mathcal S)=L_G(\mathcal S)L_G(\mathcal S)^\T+G_{\rm env}(\mathcal S).
\]
A neural model may predict these factors around an analytic/low-order baseline.

\subsection{Decoder construction}

The continuous temperature decoder is represented as
\[
 T(\mathbf r,t)-T_{\rm ref}
 =
 \sum_{m=1}^{R_T}z_m(t)\Phi_m(\mathbf r;\mathcal S).
\]
The basis functions combine:
\begin{enumerate}
\item object-local analytic Green responses of loss channels;
\item low-order interface response modes;
\item a learned continuous residual constrained to reproduce the reduced
      observation operator.
\end{enumerate}

\subsection{Nonlinear update}

For temperature-dependent material laws, a time step uses a small number of
Picard/Newton updates in reduced coordinates:
\[
 C_T(z^{(k)})\dot z^{(k+1)}
 =
 -G_T(z^{(k)})z^{(k+1)}
 +B_T(z^{(k)})p^{(k)}.
\]
The update is accepted only if the electrothermal energy residual and state
increment meet tolerances. If reduced nonlinearity cannot represent the local
constitutive change, CERTIFIED mode escalates to a local teacher correction.

\subsection{Steady-state consistency}

The steady state satisfies
\[
 G_T(z_\ast)z_\ast
 =
 B_T(z_\ast)p_\ast+b_{\rm env}.
\]
A transient integrator and steady solver use the same reduced operators. A model
that matches trajectories but fails this equation is not considered thermally
certified.
